\scanpage{16}
\mspage{15}

application canonique
\[
(*)\qquad \operatorname{Riv}(A,X)\longrightarrow\operatorname{Comp}(\widetilde A)
\]
qui est \inserted{injective, et} \uline{bijective} si $X$ paracompact.
Donc dans ce cas, choisir une rive, c'est choisir un composant de $\widetilde A$.

\uline{NB} L'ens des \deleted{composants}\inserted{rives} de \deleted{$X$} $A$ dans $X$ est une $\mathbb F_2$-algèbre \inserted{comm.} idempotente ($E^2=E$: $\forall E\in\operatorname{Riv}$) avec unité, l'application précédente est un hom. d'alg. unitaires.

Une \uline{birive} \inserted{de $A$ dans $X$} est une \uline{paire} \inserted{$\{\rho,\rho'\}$} d'éléments de $\operatorname{Riv}(A)$, telle que
\[
\rho+\rho'=1
\]
\margin{\uline{NB} Cette algèbre $\operatorname{Riv}$ est nulle ssi $A$ est ouvert, \unsure{en $=$ temps que} fermé.}
\margin{(c'est-à-dire une ``décomposition en 2 composants de $X\setminus A$ au vois. de $A$''}

\deleted{\unsure{Ceci doit être}} (ce qui implique, si $A$ n'est pas ouvert, que $\rho\ne\rho'$).
La birive est triviale si $\{\rho,\rho'\}=\{0,1\}$. On dit que $A$ disconnect $X$ loc.\textsuperscript{t} si $\exists$ birive non triviale.

\deleted{\unsure{Considérons}} Une birive définit $\{\widetilde A_1,\widetilde A_2\}$ telle que $\widetilde A=\widetilde A_1\amalg\widetilde A_2$.
On dit que \deleted{c'est une \unsure{b}} $\{A,\deleted{\{\rho,\rho'\}}\}$

\rotatebox{90}{\reflectbox{\unsure{EXTRA STRONG}}}
\rotatebox{90}{\reflectbox{\unsure{EXTRA STRONG}}}
\mspage{15}
